\section*{一、预习报告}

\subsection*{1. 实验综述}

通过预习，我了解到密立根（R.~A.~Millikan）利用电场和重力场中运动的带电油滴，测定了基本电荷 $e$，并发现所有带电体的电荷量都是 $e$ 的整数倍，从而证明了电荷的不连续性（量子性）。

本实验采用静态平衡法。油滴喷入间距为 $d$ 的两块水平平行板之间时，因摩擦而带电，所受重力为 $mg$，静电力为 $qE$。调节板间电压 $U$，使油滴静止，则

\[
mg=qE=q\frac{U}{d}.
\]

由于油滴质量很小（约 $10^{-15}\,\mathrm{kg}$），需要用匀速下降法间接测定。油滴在表面张力作用下呈小球状，$m=\frac43\pi r^3\rho$。平行板不加电压时，油滴下落一段距离后，空气黏滞阻力与重力平衡而匀速下降。由斯托克斯定律 $6\pi r\eta v=mg$（忽略空气浮力）得

\[
r=\sqrt{\frac{9\eta v}{2\rho g}}.
\]

对于半径小于 $10^{-6}\,\mathrm{m}$ 的小油滴，空气不能再看作连续介质，黏滞阻力需修正为 $6\pi r\eta v/(1+\frac{b}{pr})$，于是

\[
r=\sqrt{\frac{9\eta v}{2\rho g}\cdot\frac{1}{1+\frac{b}{pr}}}.
\]

修正项中的 $r$ 不需十分精确，可用未修正的结果代入。匀速下降的速度由下降距离 $l$ 和所用时间 $t$ 求出，即 $v=l/t$。将上述各式代入平衡条件，整理得

\[
q=\frac{18\pi}{\sqrt{2\rho g}}\left[\frac{\eta l}{t\left(1+\frac{b}{pr}\right)}\right]^{3/2}\frac{d}{U}.
\]

教材给出的常量为：$g=9.79\,\mathrm{m/s^2}$，$\eta=1.82\times10^{-5}\,\mathrm{N\cdot s/m^2}$，$\rho=981\,\mathrm{kg/m^3}$，$p=1.013\times10^5\,\mathrm{Pa}$，修正常数 $b=6.17\times10^{-6}\,\mathrm{m\cdot cmHg}$，$d=5\,\mathrm{mm}$。注意 $b/(pr)$ 中 $p$ 须以 $\mathrm{cmHg}$ 为单位，即 $p\approx76.0\,\mathrm{cmHg}$。

实验时，我将先调节测量室底座的三个螺钉，使测量室水平，保证电场方向与重力方向平行。接着用喷雾器喷入油滴，调节显微镜焦距，在显示屏上选择合适的油滴：按课件要求，应选带电量小于 10 个基本电荷、平衡电压约 $200\,\mathrm{V}$、下落时间在 $10\,\mathrm{s/mm}$ 以内且清晰可见的油滴。随后练习控制油滴，使它能在视场中上下往复运动。

正式测量时，我将把 K2 拨到“平衡”挡，调节平衡电压使油滴静止，记录 $U$；再拨到“提升”挡，把油滴移到视场上端的刻度线；然后拨到“0V”挡，在油滴匀速下落的一段内计时，记录时间 $t$ 和下落距离 $l$。我计划对同一油滴重复测量，也可换用不同油滴，共记录 10 组数据，再由上式算出各油滴的电荷量 $q_i$。

数据处理将按教材的逐次相减法进行：把 $q_i$ 从小到大排列，求相邻差值 $\Delta q=q_{i+1}-q_i$，由此先估计基本电荷 $e_0$；再由 $q_i/e_0$ 取整得到电荷数 $n_i$，计算 $e_i=q_i/n_i$，最后取平均得到 $e$。课件还给出了平均值法和多重差值法，可用来对照。课件要求求出不确定度，由于最终结果是对多个 $e_i$ 求平均，A 类分量将取 $e_i$ 平均值的标准偏差；B 类分量须以现场仪器给出的允许误差限为准，本次预习不预设数值，也不把显示分辨力当作允许误差限。

\subsection*{2. 实验目的}

\begin{enumerate}
\item 了解密立根油滴仪的结构和工作特性，学会选择和控制油滴。
\item 用静态平衡法测定电子的基本电荷量 $e$，验证电荷的不连续性。
\item 掌握逐次相减法（公约数法）等密立根油滴实验的数据处理方法。
\end{enumerate}

\subsection*{3. 实验重点与难点}

\begin{enumerate}
\item 选择合适的油滴是实验成败的关键。油滴太大，下落太快，不易测准时间，带电量也偏多，相邻电荷数难以区分；油滴太小，则受布朗运动影响明显，下落速度涨落很大。
\item 平衡电压要在油滴确实静止时读取；计时应选在油滴匀速下落的一段，避开起始加速段，起止时刻对准刻度线。跟踪测量时要随时微调焦距，防止油滴变模糊或漂出视场而丢失。
\item 测量前要调平测量室；喷雾适量，避免视场中油滴过多难以辨认。计算时注意 $b/(pr)$ 的单位，并统一换算为国际单位制。
\end{enumerate}
